Understanding whether alignment resistance in large language models (LLMs) emerges as a sudden phase transition or scales gradually with model size is critical for AI safety. We conduct a systematic behavioral study across five commercial LLMs spanning approximately 1B to 400B parameters, using three complementary evaluation protocols: jailbreak resistance under multiple attack templates, few-shot alignment erosion, and instruction conflict resolution. We apply both continuous and discrete evaluation metrics following the framework of \citet{schaeffer2023emergent} to distinguish genuine scaling phenomena from measurement artifacts. Our results show no evidence of a phase transition in alignment resistance. Instead, we observe a weak, gradual decrease in safety under hypothetical-framing attacks as model scale increases (Pearson $r = -0.765$, $p = 0.132$), complete immunity to few-shot alignment erosion at all scales, and a context-dependent vulnerability where the largest model is most susceptible to educational framing in conflict scenarios (safety score 2.7/10). Both continuous and discrete metrics track consistently, indicating these trends are not measurement artifacts. Our findings suggest that alignment resistance in well-trained commercial models degrades predictably rather than failing catastrophically, though the increasing susceptibility of larger models to sophisticated contextual framing warrants attention as models continue to scale.
